\documentclass[11pt,reqno]{amsart}

\usepackage[T1]{fontenc}
\usepackage[utf8]{inputenc}
\usepackage[margin=1in]{geometry}
\usepackage{amsmath,amssymb,amsthm,amsfonts,mathtools}
\usepackage{xcolor}
%\usepackage{microtype}
\usepackage[colorlinks=true,linkcolor=blue!70!black,citecolor=blue!70!black,urlcolor=blue!70!black]{hyperref}

\numberwithin{equation}{section}

\theoremstyle{plain}
\newtheorem{theorem}{Theorem}[section]
\newtheorem{lemma}[theorem]{Lemma}
\newtheorem{proposition}[theorem]{Proposition}
\newtheorem{corollary}[theorem]{Corollary}
\newtheorem{conjecture}[theorem]{Conjecture}
\newtheorem{problem}[theorem]{Problem}

\theoremstyle{definition}
\newtheorem{definition}[theorem]{Definition}
\newtheorem{example}[theorem]{Example}
\newtheorem{remark}[theorem]{Remark}

\newcommand{\N}{\mathbb N}
\newcommand{\R}{\mathbb R}
\newcommand{\eps}{\varepsilon}
\newcommand{\floor}[1]{\left\lfloor #1\right\rfloor}
\newcommand{\ceil}[1]{\left\lceil #1\right\rceil}
\newcommand{\pp}{\mathcal P}
\newcommand{\bb}{\mathcal B}
\newcommand{\cc}{\mathcal C}
\DeclareMathOperator{\supp}{supp}

\begin{document}

\title[An upper bound for gaps in multiplicative Sidon sets]
{A note on the bound $g(n)\ll_\varepsilon n^{10/33+\varepsilon}$}

\author{ChatGPT}
\date{}

\begin{abstract}
For a positive integer \(n\), let \(g(n)\) denote the least scale on which a
multiplicative Sidon subset of \([n]\) can be made to meet every interval.
Using the currently available exceptional-set estimate for the prime number
theorem in short intervals due to Gafni and Tao, we prove
\[
        g(n) \ll n^{10/33+o(1)}.
\]
The proof uses primes to cover the prime-rich blocks and then applies a local
lemma argument to fill the remaining prime-poor blocks.  Under the Riemann
hypothesis, the same method gives
\[
        g(n)\ll n^{1/4+o(1)}.
\]
\end{abstract}

\maketitle

\section{Introduction and external input}

Throughout the paper, we write $\N:=\{1,2,3,\dots\}$. A set $A\subseteq\N$ is
called a \emph{multiplicative Sidon set} if all products $aa'$ with $a,a'\in A$
and $a\le a'$ are distinct.

Multiplicative Sidon sets have been studied since the work of
Erd\H{o}s~\cite{Erdos1938,Erdos1968}.  In particular, the maximum size of a
multiplicative Sidon subset of \([n]:=\{1,\dots,n\}\) is
\[
        \pi(n)+\Theta\!\left(\frac{n^{3/4}}{(\log n)^{3/2}}\right).
\]
Thus the extremal theory of multiplicative Sidon sets is quite well understood
from the point of view of cardinality.  The corresponding spacing problem,
however, is much less developed.

In his 2001 list of unsolved problems,
S\'ark\"ozy~\cite[Problem~34]{Sarkozy2001} asked the following problem.

\begin{problem}[S\'ark\"ozy~\cite{Sarkozy2001}]\label{prob:34}
How small can one make the maximal gap between the consecutive elements of a
multiplicative Sidon set selected from $\{1,2,\dots,n\}$? Does there exist for
all $n$ a multiplicative Sidon set $A\subseteq\{1,2,\dots,n\}$ so that
$A\cap [m,m+\sqrt n]\neq\varnothing$ for all $1\le m<m+\sqrt n\le n$?
\end{problem}

S\'ark\"ozy observed that a positive answer to the second question would follow,
for example, from a prime-gap estimate of the shape
\(p_{i+1}-p_i=o(p_i^{1/2})\).  Such a bound is far beyond what is currently
known.

To encode Problem~\ref{prob:34} quantitatively, we use the following function.

\begin{definition}
For $n\in\N$, define $g(n)$ to be the infimum of all real numbers $L\ge 0$
such that there exists a multiplicative Sidon set
\(A\subseteq[n]:=\{1,2,\dots,n\}\) with the property that
\[
        A\cap[x,x+L]\neq\varnothing
        \qquad\text{for every }x\in\mathbb R
        \text{ with }[x,x+L]\subseteq[1,n].
\]
\end{definition}

A natural upper-bound benchmark is provided by the primes: since the set of
primes is multiplicatively Sidon, any result on primes in short intervals gives
an upper bound for \(g(n)\).  For example, the theorem of Baker, Harman and
Pintz~\cite{BakerHarmanPintz2001} implies
\[
        g(n)\ll n^{21/40}.
\]
In the opposite direction, the cardinality bound above, together with the prime
number theorem and the pigeonhole principle, gives
\[
        g(n)\ge (1+o(1))\log n .
\]

Recently, van Doorn, Monticone and Tang~\cite{DMT2026} gave an affirmative
answer to S\'ark\"ozy's explicit \(\sqrt n\)-question and proved the
power-saving bound
\[
        g(n)\ll_\varepsilon n^{\rho+\varepsilon},
        \qquad
        \rho:=\frac{13-\sqrt{69}}{10}<0.47 .
\]
The purpose of the present note is to show that the exceptional-set estimates
for primes in short intervals due to Gafni and Tao lead to the stronger bound
\[
        g(n)\ll_\varepsilon n^{10/33+\varepsilon}.
\]



We shall use the following consequence of the work of Gafni and Tao on the
exceptional set in the prime number theorem in short intervals.

\begin{theorem}[Gafni--Tao, consequence of the exceptional-set bound]
\label{thm:GT-input}
Let \(10/33<\theta<1\). For every $\eta>0$ and all sufficiently large $X$,
\[
\left|\left\{x\in[X,2X]:
\left|\sum_{x<m\le x+x^\theta}\Lambda(m)-x^\theta\right|
        \ge \frac12 x^\theta\right\}\right|
        \ll_{\theta,\eta} X^{10/11+\eta},
\]
where $|\cdot|$ denotes Lebesgue measure.
\end{theorem}

\begin{remark}[Verification of Theorem~\ref{thm:GT-input}]
We spell out the finite calculation needed from Gafni--Tao.  Let
\[
        \theta_0:=\frac{10}{33},
        \qquad
        c_0:=\frac{1}{1-\theta_0}=\frac{33}{23}.
\]
Gafni--Tao prove that, for every fixed \(0<\theta<1\),
\[
\mu(\theta)\le
\inf_{\xi>0}\ 
\sup_{\substack{0\le \sigma<1\\
        A(\sigma)\ge (1-\theta)^{-1}-\xi}}
\left((1-\theta)(1-\sigma)A(\sigma)+2\sigma-1\right).
\]
Let \(\widetilde A(\sigma)\) be the piecewise upper bound for \(A(\sigma)\)
given in Table~1 of Gafni--Tao, together with
\[
        \widetilde A(\sigma)=\frac{1}{1-\sigma}
        \qquad (0\le \sigma\le 1/2).
\]
Set
\[
        F(\sigma):=
        \frac{23}{33}(1-\sigma)\widetilde A(\sigma)+2\sigma-1.
\]
It suffices to verify that
\[
        F(\sigma)\le \frac{10}{11}
        \qquad
        \text{whenever }\widetilde A(\sigma)\ge \frac{33}{23}.
\]
The condition \(\widetilde A(\sigma)\ge 33/23\) restricts the last relevant
Table~1 entry
\[
        \widetilde A(\sigma)=\frac{3}{10\sigma-7}
        \qquad
        \left(\frac9{10}<\sigma\le \frac{31}{34}\right)
\]
to \(\sigma\le 10/11\).  No later Table~1 entry contributes.  Taking closures
of the relevant ranges, which can only enlarge the suprema, direct
differentiation of the corresponding rational functions gives the following
finite check:
\[
\begin{array}{c|c|c}
\text{range of }\sigma
&
\widetilde A(\sigma)
&
\displaystyle \sup F(\sigma)
\\ \hline
\left[\dfrac{10}{33},\dfrac12\right]
&
\dfrac{1}{1-\sigma}
&
\dfrac{23}{33}
\\[6pt]
\left[\dfrac12,\dfrac{7}{10}\right]
&
\dfrac{3}{2-\sigma}
&
\dfrac{631}{715}
\\[6pt]
\left[\dfrac{7}{10},\dfrac{19}{25}\right]
&
\dfrac{15}{3+5\sigma}
&
\dfrac{4156}{4675}
\\[6pt]
\left[\dfrac{19}{25},\dfrac{127}{167}\right]
&
\dfrac{9}{8\sigma-2}
&
\dfrac{4156}{4675}
\\[6pt]
\left[\dfrac{127}{167},\dfrac{13}{17}\right]
&
\dfrac{15}{13\sigma-3}
&
\dfrac{1625}{1837}
\\[6pt]
\left[\dfrac{13}{17},\dfrac{17}{22}\right]
&
\dfrac{6}{5\sigma-1}
&
\dfrac{985}{1122}
\\[6pt]
\left[\dfrac{17}{22},\dfrac{41}{53}\right]
&
\dfrac{2}{9\sigma-6}
&
\dfrac{608}{693}
\\[6pt]
\left[\dfrac{41}{53},\dfrac79\right]
&
\dfrac{9}{7\sigma-1}
&
\dfrac{1721}{1980}
\\[6pt]
\left[\dfrac79,\dfrac{1867}{2347}\right]
&
\dfrac{9}{8(2\sigma-1)}
&
\dfrac{1721}{1980}
\\[6pt]
\left[\dfrac{1867}{2347},\dfrac78\right]
&
\dfrac{3}{2\sigma}
&
\dfrac{277}{308}
\\[6pt]
\left[\dfrac78,\dfrac{279}{314}\right]
&
\dfrac{3}{10\sigma-7}
&
\dfrac{920849}{1022384}
\\[6pt]
\left[\dfrac{279}{314},\dfrac9{10}\right]
&
\dfrac{24}{30\sigma-11}
&
\dfrac{199}{220}
\\[6pt]
\left[\dfrac9{10},\dfrac{10}{11}\right]
&
\dfrac{3}{10\sigma-7}
&
\dfrac{10}{11}.
\end{array}
\]
Hence
\[
        F(\sigma)\le \frac{10}{11}
        \qquad
        \text{on the range }\widetilde A(\sigma)\ge \frac{33}{23}.
\]

Now fix \(\theta>10/33\), and put \(c_\theta=(1-\theta)^{-1}\).  Since
\(c_\theta>c_0\), choose \(0<\xi<c_\theta-c_0\).  If
\[
        A(\sigma)\ge c_\theta-\xi,
\]
then
\[
        \widetilde A(\sigma)\ge A(\sigma)>c_0.
\]
Using \(A(\sigma)\le \widetilde A(\sigma)\) and
\(1-\theta<1-\theta_0=23/33\), we obtain
\[
\begin{aligned}
(1-\theta)(1-\sigma)A(\sigma)+2\sigma-1
&\le
\frac{23}{33}(1-\sigma)\widetilde A(\sigma)+2\sigma-1  \\
&\le \frac{10}{11}.
\end{aligned}
\]
Therefore the Gafni--Tao bound gives
\[
        \mu(\theta)\le \frac{10}{11}
        \qquad
        (\theta>10/33).
\]
Taking \(\delta=1/2\) in the definition of the exceptional-set exponent gives
Theorem~\ref{thm:GT-input}.
\end{remark}

We shall also use the following conditional consequence of the same
exceptional-set framework.

\begin{theorem}
\label{thm:GT-RH-input}
Assume the Riemann hypothesis.  Let $0<\theta\le 1/2$.  For every $\eta>0$ and
all sufficiently large $X$,
\[
\left|\left\{x\in[X,2X]:
\left|\sum_{x<m\le x+x^\theta}\Lambda(m)-x^\theta\right|
        \ge \frac12 x^\theta\right\}\right|
        \ll_{\theta,\eta} X^{1-\theta+\eta},
\]
where $|\cdot|$ denotes Lebesgue measure.
\end{theorem}

We also use a standard elementary estimate for prime powers in very short
intervals.

\begin{lemma}[Prime powers are negligible below the square-root scale]
\label{lem:prime-powers}
Fix $0<\theta<1/2$.  Uniformly for $x\ge2$,
\[
        \sum_{\substack{x<m\le x+x^\theta\\ m=p^k,\ k\ge2}} \Lambda(m)
        \ll_\theta (\log x)^2.
\]
\end{lemma}

\begin{proof}
For a fixed $k\ge2$, the number of prime powers $p^k\in(x,x+x^\theta]$ is at
most
\[
        (x+x^\theta)^{1/k}-x^{1/k}+1
        \ll x^{\theta+1/k-1}+1.
\]
Summing over $2\le k\le \log_2(2x)$ gives
\[
        \ll x^{\theta-1/2}\log x+\log x\ll_\theta \log x
\]
prime powers.  Each contributes at most $O(\log x)$ to $\Lambda$, and the claim
follows.
\end{proof}

\section{Two combinatorial lemmas}

The first lemma is a convenient local-lemma form tailored to our selection
problem.

\begin{lemma}[A variable local lemma]
\label{lem:variable-LLL}
Let $X_1,\ldots,X_N$ be mutually independent random variables, and let
$\mathcal E$ be a finite family of bad events.  Suppose that every
$E\in\mathcal E$ is determined by at least one and at most $r$ of the variables,
where $r\ge 1$.  For a variable $X_i$, write $i\in\supp(E)$ if $E$ depends on
$X_i$.  If
\[
        \sum_{E\in\mathcal E:\ i\in\supp(E)} \mathbb P(E) \le \frac1{8r}
        \qquad (1\le i\le N),
\]
then with positive probability none of the events in $\mathcal E$ occurs.
\end{lemma}

\begin{proof}
We use the asymmetric Lov\'asz local lemma.  Put
\[
        x_E:=2\mathbb P(E) \qquad (E\in\mathcal E).
\]
Since every event has non-empty support, for any \(E\in\mathcal E\) and any
\(i\in\supp(E)\) we have
\[
        \mathbb P(E)
        \le
        \sum_{F:\ i\in\supp(F)}\mathbb P(F)
        \le \frac1{8r}.
\]
Thus \(0\le x_E\le 1/(4r)<1\).

For the dependency graph, join two distinct events when they share at least one
underlying variable.  For any fixed bad event \(E\),
\[
        \sum_{F\sim E} x_F
        \le 2\sum_{i\in\supp(E)}
              \sum_{F:\ i\in\supp(F)} \mathbb P(F)
        \le 2r\cdot \frac1{8r}=\frac14.
\]
Therefore
\[
        \prod_{F\sim E}(1-x_F)
        \ge 1-\sum_{F\sim E}x_F
        \ge \frac34,
\]
and hence
\[
        x_E\prod_{F\sim E}(1-x_F)
        \ge \frac32\mathbb P(E)
        \ge \mathbb P(E).
\]
The asymmetric Lov\'asz local lemma now gives a positive probability of avoiding
all bad events.
\end{proof}
The second lemma constructs a multiplicative Sidon transversal from a sparse
family of blocks.

\begin{lemma}[Sidon transversal for sparse blocks]
\label{lem:sparse-transversal}
Let \(H=H(n)\to\infty\), let \(\mathcal B\) be a family of pairwise disjoint
intervals of integers contained in \([n]\), each of length \(H\), and put
\(M:=|\mathcal B|\).  For each \(B\in\mathcal B\), let
\(\mathcal C_B\subseteq B\) satisfy
\[
        |\mathcal C_B|\ge \gamma H
\]
for some fixed \(\gamma>0\).  Assume that there is a fixed \(\alpha>0\) such
that
\[
        \frac{M n^\alpha}{H^2}=o(1),
        \qquad
        \frac{n^\alpha}{H^2}=o(1),
        \qquad\text{and}\qquad
        \frac{\sqrt n}{H^2}=o(1).
\]
Then there exist choices \(c_B\in\mathcal C_B\), one for each
\(B\in\mathcal B\), such that
\[
        C:=\{c_B:B\in\mathcal B\}
\]
is a multiplicative Sidon set.
\end{lemma}

\begin{proof}
Choose independently and uniformly
\[
        X_B\in\mathcal C_B \qquad (B\in\mathcal B).
\]
We shall avoid atomic bad events encoding non-trivial product coincidences.

For distinct blocks \(B_1,B_2,B_3,B_4\) and values
\(a_i\in\mathcal C_{B_i}\), include the atomic event
\[
        X_{B_i}=a_i \qquad (1\le i\le 4)
\]
whenever
\[
        a_1a_2=a_3a_4.
\]
For distinct blocks \(B_1,B_2,B_3\) and values
\(a_i\in\mathcal C_{B_i}\), include the atomic event
\[
        X_{B_1}=a_1,\qquad X_{B_2}=a_2,\qquad X_{B_3}=a_3
\]
whenever either
\[
        a_1^2=a_2a_3
        \qquad\text{or}\qquad
        a_1a_2=a_3^2.
\]
Avoiding all these atomic events is sufficient.  Indeed, consider any equality
\[
        c_{B_1}c_{B_2}=c_{B_3}c_{B_4}
\]
among selected elements.  If a block index occurs on both sides, then
cancelling the corresponding positive integer gives either the same unordered
pair of block indices, or an equality between selected values from two distinct
blocks; the latter is impossible since the blocks are disjoint.  We may
therefore assume that no block index occurs on both sides.

If the four block indices are distinct, the equality is excluded by one of the
four-variable atoms.  If exactly one side has a repeated block index, the
equality is excluded by one of the square-type atoms.  Finally, if both sides
have repeated block indices, then the equality has the form
\(c_B^2=c_{B'}^2\) with \(B\ne B'\), which again forces \(c_B=c_{B'}\) and is
impossible because the blocks are disjoint.  Thus every non-trivial product
coincidence is excluded.

We estimate, for a fixed block \(B\in\mathcal B\), the total probability of all
atomic bad events involving \(X_B\).  Let
\[
        U:=\bigcup_{B\in\mathcal B}\mathcal C_B,
        \qquad |U|\le MH.
\]
We use the divisor bound in the form
\[
        \tau(m)\ll_\alpha m^{\alpha/2}\ll_\alpha n^\alpha
        \qquad (m\le n^2).
\]

First consider four-variable atoms.  Up to an absolute factor accounting for the
possible positions of \(X_B\), fix a value \(a\in\mathcal C_B\).  After choosing
another value \(b\in U\), the number of ordered pairs \((c,d)\) with
\(cd=ab\) is at most \(\tau(ab)\).  Since the blocks are disjoint, each value
of \(U\) determines its block.  Hence the number of four-variable atoms
involving \(X_B\) is
\[
        \ll \sum_{a\in\mathcal C_B}\sum_{b\in U}\tau(ab)
        \ll_\alpha M H^2 n^\alpha .
\]
Each such atom has probability at most \((\gamma H)^{-4}\), so the total
four-variable contribution is
\[
        \ll_{\gamma,\alpha} \frac{M n^\alpha}{H^2}=o(1).
\]

Next consider square-type atoms. The two estimates below account, up to an absolute factor, for all possible
positions of the fixed variable \(X_B\) in a three-variable atom. For atoms of the shape \(a^2=bc\) with
\(a\in\mathcal C_B\), the number of possible ordered pairs \((b,c)\) is at most
\[
        \tau(a^2)\ll_\alpha n^\alpha.
\]
Therefore the total contribution of these atoms is
\[
        \ll_{\gamma,\alpha} \frac{H n^\alpha}{H^3}
        =
        \frac{n^\alpha}{H^2}
        =
        o(1).
\]
For atoms of the shape \(ab=c^2\) with \(a\in\mathcal C_B\), write uniquely
\[
        a=rs^2
\]
with \(r\) squarefree.  Then every positive integer solution has
\[
        b=rt^2,
        \qquad c=rst
\]
for some positive integer \(t\).  Since \(b\le n\), there are at most
\(O(\sqrt n)\) possible \(t\).  Hence the number of such atoms involving
\(X_B\) is \(O(H\sqrt n)\), and their total contribution is
\[
        \ll_\gamma \frac{H\sqrt n}{H^3}
        =
        \frac{\sqrt n}{H^2}
        =
        o(1).
\]

Combining the three estimates, the total probability of all atomic bad events
involving any fixed variable is \(o(1)\).  Since every atomic event depends on
at most four variables, this total is at most \(1/(8\cdot 4)\) for all
sufficiently large \(n\).  Lemma~\ref{lem:variable-LLL} therefore applies and
gives a choice avoiding all atomic bad events.  For this choice, the selected
set \(C\) is multiplicative Sidon.
\end{proof}

\section{Counting prime-poor blocks}

Fix \(0<\eps<1/10\) and set
\[
        \kappa:=\frac{10}{33}+\eps,
        \qquad H:=\ceil{n^\kappa}.
\]
We shall eventually prove $g(n)\le 2H$ for all sufficiently large $n$.
Choose a number $\theta$ with
\[
        \frac{10}{33}<\theta<\kappa<\frac12.
\]
For instance one may take $\theta=10/33+\eps/2$.
Let
\[
        B_r:=\{(r-1)H+1,(r-1)H+2,\ldots,rH\},
        \qquad 1\le r\le t:=\floor{n/H}.
\]
These are the full blocks of length $H$.

Fix once and for all a large integer $R\ge100$.
A block $B_r$ is called \emph{bad} if either
\[
        B_r\cap [1,2\sqrt n]\ne\varnothing
\]
or $B_r$ contains at most $R$ primes.  All other full blocks are called
\emph{good}.  Thus every good block lies above $2\sqrt n$ and contains at least
$R+1$ primes.

\begin{lemma}[Number of bad blocks]
\label{lem:bad-blocks}
Let $M$ be the number of bad blocks.  Then
\[
M \ll \frac{\sqrt n}{H}+1+\frac{n^{10/11+o(1)}}{H}.
\]
In particular, if $H=n^{10/33+\eps+o(1)}$, then
\[
        \frac{M n^{o(1)}}{H^2}=o(1).
\]
\end{lemma}

\begin{proof}
There are $O(\sqrt n/H+1)$ bad blocks meeting $[1,2\sqrt n]$.  We now count the
remaining bad blocks, which are disjoint from $[1,2\sqrt n]$.

Let \(B=[u+1,u+H]\) be one of these remaining bad blocks.  Then
\(u+1>2\sqrt n\), and \(B\) contains at most \(R\) primes.  Since
\(\theta<\kappa\), we have \(n^\theta=o(H)\).  Hence, for all sufficiently
large \(n\) and all
\[
        x\in [u+1,u+H/2],
\]
every integer \(m\) with
\[
        x<m\le x+x^\theta
\]
lies in \(B\).  Indeed, \(x+x^\theta\le u+H/2+n^\theta\le u+H\) for all large
\(n\).  Therefore this short interval contains at most \(R\) primes.  By
Lemma~\ref{lem:prime-powers}, the contribution of prime powers to
\(\sum_{x<m\le x+x^\theta}\Lambda(m)\) is \(O((\log x)^2)\), while the
contribution of primes is \(O_R(\log x)\).  Since \(x\ge u+1>2\sqrt n\), this
gives
\[
        \sum_{x<m\le x+x^\theta}\Lambda(m)=o(x^\theta)
\]
uniformly for such \(x\).  In particular, for all large \(n\),
\[
\left|\sum_{x<m\le x+x^\theta}\Lambda(m)-x^\theta\right|
        \ge \frac12 x^\theta.
\]
Thus each high bad block contributes a subinterval of \(x\)-values of measure
\(\gg H\) to the exceptional set in Theorem~\ref{thm:GT-input}.  These
subintervals are disjoint as \(B\) varies.

For all sufficiently large \(n\), all these \(x\)-values are at least
\(2\sqrt n\), so the dyadic intervals to which Theorem~\ref{thm:GT-input} is
applied have lower endpoints tending to infinity.  Summing the exceptional-set
estimate over dyadic intervals with \(X\le n\) gives total exceptional measure
\[
        \ll n^{10/11+o(1)}
\]
in \([1,n]\).  Hence the number of high bad blocks is
\[
        \ll \frac{n^{10/11+o(1)}}{H}.
\]
This proves the asserted bound for \(M\).

Finally,
\[
        \frac{M n^{o(1)}}{H^2}
        \ll
        n^{1/2-3\kappa+o(1)}
        +n^{-2\kappa+o(1)}
        +n^{10/11-3\kappa+o(1)}
        =o(1),
\]
because $\kappa>10/33>1/4$ and $10/11=3\cdot 10/33$.
\end{proof}

\section{Choosing primes in the good blocks}

Let $\mathcal G$ be the set of good blocks and $\mathcal B$ the set of bad
blocks.  From each good block $G\in\mathcal G$, choose a set $P_G$ of exactly
$R$ primes contained in $G$.  Every such prime is larger than $2\sqrt n$.
We choose one prime $q_G\in P_G$ independently and uniformly at random.

\begin{lemma}[A good choice of the good-block primes]
\label{lem:choose-good-primes}
There is a choice of primes
\[
        Q:=\{q_G:G\in\mathcal G\}
\]
such that for every bad block $B\in\mathcal B$, at least $H/2$ integers in $B$
are not divisible by any prime in $Q$.
\end{lemma}

\begin{proof}
Fix a bad block $B$.  For a good block $G$, let $Y_G(B)$ be the indicator of
the event that the randomly chosen prime $q_G$ divides at least one integer of
$B$.  The variables $Y_G(B)$ are independent as $G$ varies.  Since every
$q_G>2\sqrt n$ and every integer $a\le n$ is divisible by at most one prime
larger than $\sqrt n$, the total number of pairs
\[
        (a,q),\qquad a\in B,\quad q\in\bigcup_{G\in\mathcal G}P_G,\quad q\mid a,
\]
is at most $H$.  Consequently
\[
        \mathbb E\sum_{G\in\mathcal G}Y_G(B) \le \frac{H}{R}.
\]
Let
\[
        S_B:=\sum_{G\in\mathcal G}Y_G(B),
        \qquad
        \mu_B:=\mathbb E S_B\le \frac{H}{R}.
\]
By the standard Chernoff bound for sums of independent Bernoulli variables,
\[
        \mathbb P\left(S_B>\frac H2\right)
        \le
        \left(\frac{e\mu_B}{H/2}\right)^{H/2}
        \le
        \left(\frac{2e}{R}\right)^{H/2}
        \le \exp(-cH),
\]
with an absolute constant \(c>0\), since \(R\ge100\). Since $H=o(\sqrt n)$ in the applications
below, every chosen prime $q_G>2\sqrt n$ is larger than $H$.  Hence such a prime
divides at most one integer in any fixed block of length $H$.  Consequently, for
a fixed bad block $B$, the number of integers in $B$ divisible by at least one
chosen prime is at most
\[
        \sum_{G\in\mathcal G}Y_G(B).
\]
There are at most $n/H$ bad blocks, and $H\to\infty$ polynomially.  Hence the
union bound shows that, with positive probability, every bad block contains at
most $H/2$ integers divisible by the chosen primes.  This gives the desired
deterministic choice of $Q$.
\end{proof}

Fix such a choice of $Q$.  For every bad block $B$, define
\[
        \mathcal C_B:=\{a\in B:q\nmid a\text{ for every }q\in Q\}.
\]
By Lemma~\ref{lem:choose-good-primes},
\[
        |\mathcal C_B|\ge H/2
        \qquad (B\in\mathcal B).
\]

\section{Completion of the proof}

\begin{theorem}
\label{thm:main}
For every $\eps>0$,
\[
        g(n)\ll_\eps n^{10/33+\eps}.
\]
Equivalently,
\[
        g(n)\le n^{10/33+o(1)}.
\]
\end{theorem}

\begin{proof}
It is enough to prove the result for all sufficiently large $n$, since the
implicit constant may be enlarged to handle finitely many exceptional values.
We may also suppose $0<\eps<1/10$.  Put
\[
        \kappa=\frac{10}{33}+\eps,
        \qquad H=\ceil{n^\kappa}.
\]
Use the block decomposition above.

By Lemma~\ref{lem:bad-blocks}, if \(M\) denotes the number of bad blocks, then
\[
        M \ll \frac{\sqrt n}{H}+1+\frac{n^{10/11+o(1)}}{H}.
\]
Take \(\alpha:=\eps\).  Since \(H=n^{\kappa+o(1)}\), we have
\[
\begin{aligned}
        \frac{M n^\alpha}{H^2}
        &\ll
        n^{1/2+\alpha-3\kappa+o(1)}
        +
        n^{\alpha-2\kappa+o(1)}
        +
        n^{10/11+\alpha-3\kappa+o(1)}        \\
        &=o(1),
\end{aligned}
\]
because \(\kappa=10/33+\eps\), \(\alpha=\eps\), and \(\kappa>10/33>1/4\). Also
\[
        \frac{n^\alpha}{H^2}
        =
        n^{\alpha-2\kappa+o(1)}
        =
        o(1),
\]
and
\[
        \frac{\sqrt n}{H^2}
        =
        n^{1/2-2\kappa+o(1)}
        =
        o(1),
\]
because \(\kappa>10/33>1/4\).  Therefore Lemma~\ref{lem:sparse-transversal},
applied with \(\gamma=1/2\), gives a choice
\[
        c_B\in\mathcal C_B \qquad (B\in\mathcal B)
\]
such that
\[
        C:=\{c_B:B\in\mathcal B\}
\]
is multiplicative Sidon.

Now define
\[
        A:=Q\cup C.
\]
Every good block is met by $Q$, and every bad block is met by $C$.  Hence
$A$ meets every full block $B_r$, $1\le r\le t$.

We next check that $A$ is multiplicative Sidon.  Suppose
\[
        a_1a_2=a_3a_4,
        \qquad a_i\in A.
\]
For each $q\in Q$, compare the $q$-adic valuations of the two sides.  The primes
in $Q$ are distinct, and no element of $C$ is divisible by any prime in $Q$ by
construction.  Hence the multiplicity with which $q$ occurs among
$\{a_1,a_2\}$ is equal to the multiplicity with which it occurs among
$\{a_3,a_4\}$.  Therefore the multiset of $Q$-elements appearing on the left is
the same as the multiset of $Q$-elements appearing on the right.  Cancelling this common multiset, we are left with an equality involving only
elements of \(C\), with the same number of remaining factors on the two sides.
Thus the remaining equality is either \(1=1\), an equality of single
\(C\)-elements, or a product equality between two \(C\)-elements on each side.
In the first two cases it is plainly trivial, and in the last case it is
trivial because \(C\) is multiplicative Sidon.  Thus the original equality is
also trivial, and \(A\) is multiplicative Sidon.

It remains to convert block coverage into interval coverage.  Let
$[x,x+2H]\subseteq[1,n]$.  Put
\[
        r:=\floor{\frac{x-1}{H}}+2.
\]
Then
\[
        x < (r-1)H+1\le rH < x+2H.
\]
Since \(x+2H\le n\), this gives \(rH\le n\), and hence
\[
        r\le \floor{n/H}=t.
\]
Therefore the full block \(B_r\) is contained in \([x,x+2H]\), and hence this
interval meets \(A\).

Thus every interval of length $2H$ contained in $[1,n]$ meets the multiplicative
Sidon set $A$.  Consequently
\[
        g(n)\le 2H\ll n^\kappa=n^{10/33+\eps}.
\]
This proves the theorem.
\end{proof}

\begin{remark}
The exponent $10/33$ enters only through the exceptional-set estimate
$\mu(\theta)\le10/11$ for $\theta>10/33$.  The local-lemma filling step requires
roughly $M\ll H^2$, and the exceptional-set estimate gives
$M\ll n^{10/11+o(1)}/H$.  Thus one needs
\[
        \frac{n^{10/11+o(1)}}{H^3}=o(1),
\]
which is exactly $H\gg n^{10/33+o(1)}$.
\end{remark}





\section{Conditional improvement under the Riemann hypothesis}

The same argument gives a sharper conditional estimate.

\begin{theorem}
\label{thm:RH-main}
Assume the Riemann hypothesis.  Then, for every $\eps>0$,
\[
        g(n)\ll_\eps n^{1/4+\eps}.
\]
\end{theorem}

\begin{proof}
It is enough to prove the result for all sufficiently large $n$, and we may
suppose $0<\eps<1/10$.  Put
\[
        \kappa:=\frac14+\eps,
        \qquad H:=\ceil{n^\kappa},
\]
and choose
\[
        \theta:=\frac14+\frac{\eps}{2}.
\]
Then
\[
        0<\theta<\kappa<\frac12.
\]

As before, decompose $[1,n]$ into full blocks
\[
        B_r:=\{(r-1)H+1,(r-1)H+2,\ldots,rH\},
        \qquad 1\le r\le t:=\floor{n/H}.
\]
Fix a sufficiently large integer $R$, and call a block bad if it either meets
$[1,2\sqrt n]$ or contains at most $R$ primes.  All other full blocks are called
good.

We first count the bad blocks.  The number of bad blocks meeting
$[1,2\sqrt n]$ is
\[
        O\left(\frac{\sqrt n}{H}+1\right).
\]
Now let \(B=[u+1,u+H]\) be a bad block disjoint from \([1,2\sqrt n]\).
Then \(u+1>2\sqrt n\), and \(B\) contains at most \(R\) primes.  Since
\(\theta<\kappa\), we have \(n^\theta=o(H)\).  Hence, for all sufficiently
large \(n\) and all
\[
        x\in [u+1,u+H/2],
\]
every integer \(m\) with
\[
        x<m\le x+x^\theta
\]
lies in \(B\).  Therefore this short interval contains at most \(R\) primes.
By Lemma~\ref{lem:prime-powers}, the contribution of prime powers to
\[
        \sum_{x<m\le x+x^\theta}\Lambda(m)
\]
is \(O((\log x)^2)\), while the contribution of primes is \(O_R(\log x)\).
Since \(x\ge u+1>2\sqrt n\), we get
\[
        \sum_{x<m\le x+x^\theta}\Lambda(m)=o(x^\theta)
\]
uniformly for such \(x\).  Therefore each high bad block contributes a set of
\(x\)-values of measure \(\gg H\) to the exceptional set in
Theorem~\ref{thm:GT-RH-input}.  These sets are disjoint as the block varies.

For all sufficiently large \(n\), the relevant \(x\)-values are at least
\(2\sqrt n\), so the lower endpoints in the dyadic decomposition tend to
infinity.  Summing Theorem~\ref{thm:GT-RH-input} over dyadic intervals gives
total exceptional measure
\[
        \ll n^{1-\theta+o(1)}
\]
in \([1,n]\).  Hence, if \(M\) denotes the number of bad blocks, then
\[
        M\ll \frac{\sqrt n}{H}+1+\frac{n^{1-\theta+o(1)}}{H}.
\]
Taking \(\alpha:=\eps\), the above bound for \(M\) gives
\[
\begin{aligned}
        \frac{M n^\alpha}{H^2}
        &\ll
        n^{1/2+\alpha-3\kappa+o(1)}
        +
        n^{1-\theta+\alpha-3\kappa+o(1)}      \\
        &=o(1),
\end{aligned}
\]
because
\[
        3\kappa=\frac34+3\eps,
        \qquad
        1-\theta=\frac34-\frac{\eps}{2},
        \qquad
        \alpha=\eps.
\]
Also
\[
        \frac{n^\alpha}{H^2}
        =
        n^{\alpha-2\kappa+o(1)}
        =
        o(1),
\]
and
\[
        \frac{\sqrt n}{H^2}
        =
        n^{1/2-2\kappa+o(1)}
        =
        n^{-2\eps+o(1)}
        =
        o(1).
\]

Choose primes in the good blocks exactly as in
Lemma~\ref{lem:choose-good-primes}.  The same proof applies here, since
$H=o(\sqrt n)$ and every selected good-block prime is larger than $2\sqrt n$.
Thus we obtain a set $Q$ of one prime from each good block such that, for every
bad block $B$, the set
\[
        \mathcal C_B:=\{a\in B:q\nmid a\text{ for every }q\in Q\}
\]
satisfies
\[
        |\mathcal C_B|\ge H/2.
\]
By Lemma~\ref{lem:sparse-transversal}, we may choose one element
$c_B\in\mathcal C_B$ from each bad block so that
\[
        C:=\{c_B:B\text{ bad}\}
\]
is multiplicative Sidon.

Set
\[
        A:=Q\cup C.
\]
Then $A$ meets every full block.  The same valuation argument used in the proof
of Theorem~\ref{thm:main} shows that $A$ is multiplicative Sidon.  Finally, the
same block-covering argument gives that every interval of length $2H$ contained
in $[1,n]$ contains one full block and hence meets $A$.  Therefore
\[
        g(n)\le 2H\ll n^{1/4+\eps},
\]
as claimed.
\end{proof}

\begin{remark}
The exponent $1/4$ is the natural endpoint of the present filling step.  Indeed,
even if the exceptional-set estimate were stronger, Lemma~\ref{lem:sparse-transversal}
still requires
\[
        \frac{\sqrt n}{H^2}=o(1),
\]
which comes from square-type product coincidences of the form $ab=c^2$.
Thus this version of the argument cannot by itself go below
$n^{1/4+o(1)}$.
\end{remark}



















\begin{thebibliography}{99}

\bibitem{AlonSpencer}
N.~Alon and J.~H.~Spencer,
\emph{The Probabilistic Method},
4th edition, Wiley, 2016.

\bibitem{BakerHarmanPintz2001}
R.~C. Baker, G.~Harman and J.~Pintz,
\emph{The difference between consecutive primes, II},
Proc. London Math. Soc. (3) \textbf{83} (2001), no.~3, 532--562.


\bibitem{DMT2026}
W.~van Doorn, P.~Monticone and Q.~Tang,
\emph{Gaps in Multiplicative Sidon Sets},
arXiv:2605.02064, 2026.
\url{https://arxiv.org/abs/2605.02064}

\bibitem{Erdos1938}
P.~Erd\H{o}s,
\emph{On sequences of integers no one of which divides the product of two others
and on some related problems},
Mitt. Forsch.-Inst. Math. Mech. Univ. Tomsk \textbf{2} (1938), 74--82.

\bibitem{Erdos1968}
P.~Erd\H{o}s,
\emph{On some applications of graph theory to number theoretic problems},
Publ. Ramanujan Inst. No.~1 (1968), 131--136.

\bibitem{GafniTao2025}
A.~Gafni and T.~Tao,
\emph{On the number of exceptional intervals to the prime number theorem in short intervals},
arXiv:2505.24017, 2025.
\url{https://arxiv.org/abs/2505.24017}

\bibitem{Sarkozy2001}
A.~S\'ark\"ozy,
\emph{Unsolved Problems in Number Theory},
Period. Math. Hungar. \textbf{42} (2001), no.~1--2, 17--35.

\end{thebibliography}

\end{document}
